\section{Обучение ранжирующей модели и оффлайн-метрики}

\subsection{Модель}

Итоговая модель --- \texttt{CatBoostRanker} с listwise-лоссом YetiRank~\cite{yetirank}, оптимизирующим порядок объектов внутри группы; группой выступает пользователь.
Гиперпараметры вынесены в \texttt{configs/train.yaml} (таблица~\ref{tab:hyperparams}) и автоматически фиксируются в ClearML при каждом запуске.

\begin{table}[H]
\centering
\caption{Основные гиперпараметры итоговой модели}
\label{tab:hyperparams}
\small
\begin{tabularx}{\textwidth}{@{}l l X@{}}
\toprule
\textbf{Параметр} & \textbf{Значение} & \textbf{Назначение} \\
\midrule
loss\_function & YetiRank & listwise-ранжирование внутри групп \\
eval\_metric & NDCG@10 (exp) & метрика отбора лучшей итерации \\
iterations & 2000 & верхний предел деревьев \\
learning\_rate & 0.05 & темп обучения \\
depth & 8 & глубина деревьев \\
l2\_leaf\_reg & 3.0 & L2-регуляризация листьев \\
early\_stopping\_rounds & 100 & остановка при стагнации NDCG@10 \\
task\_type / devices & GPU / 0 & обучение на GPU \\
allow\_cpu\_fallback & true & автоматический переход на CPU без CUDA \\
\bottomrule
\end{tabularx}
\end{table}

Обучение итоговой модели выполнено на выделенном GPU-сервере.

\subsection{Метрики на этапе обучения}

CatBoost вычисляет метрики ранжирования на валидации прямо в процессе обучения (секция \texttt{custom\_metric}): NDCG@10, MAP@10, MRR@10, Precision@10, Recall@10.
После обучения полная история значений извлекается из \texttt{get\_evals\_result()} и по итерациям отправляется в ClearML, где кривые доступны для сравнения между экспериментами.

\subsection{Собственный модуль оффлайн-метрик}

Помимо встроенных метрик CatBoost, реализован собственный модуль \texttt{src/metrics.py}, применяемый к предсказаниям на валидации и тесте после обучения.
Все метрики групповые: значение считается по каждому пользователю и усредняется.
Реализация полностью векторизована (сортировка \texttt{lexsort} по паре <<группа, $-$скор>> и сегментные суммы \texttt{reduceat}), что позволяет обсчитывать миллионы строк без циклов по группам.

Для группы $u$ с ранжированием по убыванию скора и релевантностями $rel_i$~\cite{ndcg}:
\begin{equation}
\text{DCG@}k = \sum_{i=1}^{k} \frac{2^{rel_i} - 1}{\log_2(i+1)}, \qquad
\text{NDCG@}k = \frac{\text{DCG@}k}{\text{IDCG@}k},
\end{equation}
где $\text{IDCG@}k$ --- DCG идеального порядка.
Бинарные метрики (позитив: $\text{label} \ge 1$):
\begin{equation}
\text{AP@}k = \frac{1}{\min(P, k)} \sum_{i=1}^{k} \text{Prec@}i \cdot [rel_i \ge 1], \qquad
\text{MRR@}k = \frac{1}{r_{\text{first}}},
\end{equation}
где $P$ --- число позитивов в группе, а $r_{\text{first}}$ --- позиция первого позитива в выдаче; а также HitRate@k, Recall@k, Precision@k.
GAUC --- усредненный по пользователям ROC-AUC, вычисляемый через статистику Манна--Уитни по рангам.
Группы без позитивов (для NDCG --- с нулевым IDCG) исключаются из усреднения.

Корректность модуля подтверждена сверкой NDCG со scikit-learn~\cite{sklearn-docs} и ручными расчетами в тестах (раздел~4).

\subsection{Трекинг экспериментов}

Каждый запуск обучения создает эксперимент в проекте ClearML, содержащий:
\begin{itemize}
    \item полную конфигурацию модели и путей (секции \texttt{config}, \texttt{model});
    \item кривые всех метрик по итерациям (закладка Scalars, рисунок~\ref{fig:clearml-scalars});
    \item финальные оффлайн-метрики по обоим сплитам (Single values);
    \item артефакты: файл модели \texttt{.cbm}, таблицу важностей признаков, JSON с метриками.
\end{itemize}

\begin{figure}[H]
\centering
\includegraphics[width=\textwidth]{clearml_scalars.png}
\caption{Кривые метрик обучения по итерациям в ClearML (закладка Scalars эксперимента обучения): MAP@10, MRR@10, NDCG@10 на обучении и валидации, сводная таблица финальных значений}
\label{fig:clearml-scalars}
\end{figure}

\subsection{Результаты обучения и оценки модели}

Итоговая модель обучена на GPU-сервере по полному конфигу (лимит 2000 итераций с ранней остановкой).
Оффлайн-метрики, рассчитанные модулем \texttt{src/metrics.py} на полных валидационном и тестовом сплитах, приведены в таблице~\ref{tab:metrics}.

\begin{table}[H]
\centering
\caption{Оффлайн-метрики итоговой модели на полных сплитах}
\label{tab:metrics}
\small
\begin{tabularx}{\textwidth}{@{}l r r r r r r@{}}
\toprule
\textbf{Сплит} & \textbf{NDCG@10} & \textbf{MAP@10} & \textbf{MRR@10} & \textbf{HitRate@10} & \textbf{Recall@10} & \textbf{GAUC} \\
\midrule
validation & 0.412 & 0.387 & 0.812 & 0.987 & 0.214 & 0.812 \\
test       & 0.408 & 0.381 & 0.804 & 0.985 & 0.209 & 0.807 \\
\bottomrule
\end{tabularx}
\end{table}

Модель демонстрирует высокое качество ранжирования: GAUC~0.81 означает, что для случайной пары <<позитив--негатив>> одного пользователя модель в 81\% случаев ставит позитив выше; MRR@10~$\approx$~0.81 --- первый релевантный ролик в среднем оказывается на 1--2 позиции выдачи, а HitRate@10~$\approx$~0.99 --- почти у каждого пользователя в топ-10 есть релевантное видео.
Близость метрик на валидации и тесте (расхождение менее 1~п.\,п.) свидетельствует об отсутствии переобучения под валидацию и корректности темпорального разбиения.
Умеренные абсолютные значения Recall@10 объясняются большим числом позитивов на пользователя в оффлайн-логе: в топ-10 физически помещается лишь малая их доля.

Наиболее важными признаками по \texttt{get\_feature\_importance} оказались агрегаты релевантности: средний \texttt{label} видео (\texttt{item\_mean\_label}), средний \texttt{label} пары пользователь--автор (\texttt{ua\_mean\_label}) и средние доли досмотра видео и пользователя --- то есть признаки всех трех групп вносят вклад в ранжирование.
Кривые метрик по итерациям и артефакты обучения зафиксированы в ClearML.
